% Partie : théorie de la démonstration
% Chapitre : normalisation
% Section : conversions par réduction

\documentclass[../../../include/open-logic-section]{subfiles}

\begin{document}

\olfileid[fr]{pt}{nor}{red}

\olsection{Conversions par réduction}

Une coupure de longueur~$1$ dans une preuve~$\delta$ de
\Log{N1i} comporte une seule occurrence de formule :
elle est la conclusion d'une introduction, ou de
\FalseInt, et la prémisse majeure d'une élimination.
Une \emph{conversion par réduction} supprime
ce couple d'inférences. Des coupures peuvent alors
apparaître ou s'allonger, et des sous-preuves
peuvent être dupliquées. La diminution de la
mesure exige donc les conditions de sélection
de la coupure maximale ; elle ne résulte pas
de la seule disparition du couple.

Prenons le couple \Intro{\lif}, \Elim{\lif},
avec $!A \ident !B \lif !C$. Une sous-preuve
$\delta_1$ se termine par
\[
\AxiomC{$\Discharge{!B}{x}$}
\RightLabel{$\delta_2$}
\DeduceC{$!C$}
\DischargeRule{\Intro{\lif}}{x}
\UnaryInfC{$!B \lif !C$}
\AxiomC{}
\RightLabel{$\delta_3$}
\DeduceC{$!B$}
\RightLabel{\Elim{\lif}}
\BinaryInfC{$!C$}
\DisplayProof
\]
Nous la remplaçons par
\[
\AxiomC{}
\RightLabel{$\delta_3$}
\DeduceC{$!B$}
\RightLabel{$\Subst{\delta_2}{\delta_3}{!B^x}$}
\DeduceC{$!C$}
\DisplayProof
\]
La notation $\Subst{\delta_2}{\delta_3}{!B^x}$
désigne le greffage de toute la preuve~$\delta_3$
à chacune des occurrences d'hypothèse~$!B$
de la classe désignée~$x$ dans~$\delta_2$.
Si cette classe est vide, aucun greffage n'a lieu.

Avant la substitution, nous choisissons les
étiquettes de façon à éviter toute capture :
la classe déchargée~$x$ et les autres classes
déchargées dans la preuve receveuse sont renommées,
au besoin, pour ne pas rencontrer les hypothèses
ouvertes de la preuve insérée. Les occurrences
de la classe désignée sont ainsi supprimées ;
une hypothèse ouverte de la preuve insérée
ne se trouve pas déchargée par inadvertance.
Chaque inférence de décharge de~$\delta_2$
conserve sa classe correspondante, et chaque
copie d'une inférence de~$\delta_3$ décharge
les copies des hypothèses qui lui appartiennent.
Les constantes propres sont rendues fraîches
avant greffage et renommées séparément dans
chaque copie pour rétablir la régularité.

La conclusion reste~$!C$. Les hypothèses
ouvertes de la nouvelle preuve sont parmi
celles de~$\delta_1$, mais elles ne sont
pas nécessairement toutes conservées.
Si la décharge de~$x$ était vacante, la preuve
d'argument~$\delta_3$ disparaît ; ses
dépendances propres peuvent donc disparaître
également. La projection d'une conjonction
ou le choix d'une branche de disjonction
peut aussi supprimer des hypothèses inutiles.

Les coupures entièrement à l'extérieur de
$\delta_1$ ne sont pas modifiées, sauf celles
qui traversent la frontière du remplacement,
qu'il faut examiner séparément.
À l'intérieur des sous-preuves conservées,
chaque copie d'une coupure conserve ses
formules, son rang et sa longueur.
Cela n'affirme pas que le \emph{nombre} de
ces copies demeure le même.

Deux difficultés méritent un examen détaillé.
\begin{enumerate}
  \item Le remplacement des occurrences
  de~$!B^x$ par~$\delta_3$ peut multiplier
  les coupures de cette dernière. Si certaines
  sont maximales, la longueur de coupure
  de la nouvelle preuve peut augmenter.
  On choisit donc une coupure maximale
  supérieure et située le plus à droite,
  dans le sens précisé par le théorème :
  les sous-preuves susceptibles d'être
  copiées ne doivent pas contenir de
  coupure maximale, ni participer par
  une prémisse voisine à un segment
  maximal. Dans l'exemple d'implication,
  l'argument~$\delta_3$ se trouve à
  droite de la prémisse majeure.

  \item Le greffage peut créer une coupure
  sur~$!B$. Cela arrive si une hypothèse
  de la classe~$x$ était une prémisse
  majeure d'élimination et si~$\delta_3$
  se termine par une introduction,
  par~\FalseInt, par~\Elim{\lor} ou
  par~\Elim{\lexists}. Une occurrence
  qui n'était qu'une hypothèse devient
  alors un détour immédiat, ou la fin
  d'un segment plus long. Si~$!B^x$
  était une prémisse mineure de
  \Elim{\lor} ou de \Elim{\lexists},
  le segment déjà présent peut aussi
  s'allonger. Toutefois,
  $\depth{!B} < \depth{!B \lif !C}$.
  Ces coupures sont donc de rang
  inférieur à celui de la coupure
  supprimée. Il faut également examiner
  la frontière inférieure : une
  éventuelle nouvelle coupure sur le
  conséquent~$!C$ a elle aussi un rang
  inférieur à celui de l'implication.
\end{enumerate}

Avec ce choix de la coupure, aucune coupure
maximale n'est dupliquée et les nouvelles
coupures ont un rang inférieur. Si d'autres
coupures de l'ancien rang maximal demeurent,
la longueur de coupure diminue ; sinon,
le rang de coupure baisse.

Une conversion par réduction est aussi appelée
\emph{conversion de détour}.
Quelques schémas sont donnés dans le
\olref{tab:reduction-conversions} ;
les autres sont laissés en exercice.
Pour les schémas quantifiés, les substitutions
de termes sont effectuées sans capture.
Les constantes propres intérieures sont
choisies fraîches par rapport au terme inséré.
Dans le cas universel, la généralisation
supprimée permet de substituer le terme
fermé à sa constante propre dans la
sous-preuve. Dans le cas existentiel, on
effectue d'abord cette substitution dans
la branche mineure, puis le greffage aux
occurrences désignées de l'hypothèse ainsi
instanciée. Les conditions de fraîcheur
préservent les hypothèses restantes et
la conclusion.

\begin{table}
\begin{multline*}
\bottomAlignProof
\AxiomC{$\Discharge{!B}{x}$}
\RightLabel{$\delta_2$}
\shortDeduce
\DeduceC{$!C$}
\DischargeRule{\Intro{\lif}}{x}
\UnaryInfC{$!B \lif !C$}
\AxiomC{}
\RightLabel{$\delta_3$}
\shortDeduce
\DeduceC{$!B$}
\RightLabel{\Elim{\lif}}
\BinaryInfC{$!C$}
\DisplayProof
\quad \longrightarrow\quad
\shoveright{\bottomAlignProof
\AxiomC{}
\RightLabel{$\delta_3$}
\shortDeduce
\DeduceC{$!B$}
\RightLabel{$\Subst{\delta_2}{\delta_3}{!B^x}$}
\shortDeduce
\DeduceC{$!C$}
\DisplayProof}
\\[4ex]
\shoveleft{\bottomAlignProof
\AxiomC{}
\RightLabel{$\delta_2$}
\shortDeduce
\DeduceC{$!B$}
\AxiomC{}
\RightLabel{$\delta_3$}
\shortDeduce
\DeduceC{$!C$}
\RightLabel{\Intro{\land}}
\BinaryInfC{$!B \land !C$}
\RightLabel{\Elim{\land}[1]}
\UnaryInfC{$!B$}
\DisplayProof
\quad \longrightarrow\quad
\bottomAlignProof
\AxiomC{}
\RightLabel{$\delta_2$}
\shortDeduce
\DeduceC{$!B$}
\DisplayProof}\qquad
\shoveright{\bottomAlignProof
\AxiomC{}
\RightLabel{$\delta_2$}
\shortDeduce
\DeduceC{$!B$}
\AxiomC{}
\RightLabel{$\delta_3$}
\shortDeduce
\DeduceC{$!C$}
\RightLabel{\Intro{\land}}
\BinaryInfC{$!B \land !C$}
\RightLabel{\Elim{\land}[2]}
\UnaryInfC{$!C$}
\DisplayProof
\quad \longrightarrow\quad
\bottomAlignProof
\AxiomC{}
\RightLabel{$\delta_3$}
\shortDeduce
\DeduceC{$!C$}
\DisplayProof}
\\[4ex]
\bottomAlignProof
\AxiomC{}
\RightLabel{$\delta_2$}
\shortDeduce
\DeduceC{$!B$}
\RightLabel{\Intro{\lor}[1]}
\UnaryInfC{$!B \lor !C$}
\AxiomC{$\Discharge{!B}{x}$}
\RightLabel{$\delta_3$}
\shortDeduce
\DeduceC{$!D$}
\AxiomC{$\Discharge{!C}{y}$}
\RightLabel{$\delta_4$}
\shortDeduce
\DeduceC{$!D$}
\DischargeRule{\Elim{\lor}}{x\,y}
\TrinaryInfC{$!D$}
\DisplayProof
\quad \longrightarrow\quad
\shoveright{\bottomAlignProof
\AxiomC{}
\RightLabel{$\delta_2$}
\shortDeduce
\DeduceC{$!B$}
\RightLabel{$\Subst{\delta_3}{\delta_2}{!B^x}$}
\shortDeduce
\DeduceC{$!D$}
\DisplayProof}
\\[2ex]
\shoveleft{\bottomAlignProof
\AxiomC{}
\RightLabel{$\delta_2$}
\shortDeduce
\DeduceC{$!B(c)$}
\RightLabel{\Intro{\lforall}}
\UnaryInfC{$\lforall[x][!B(x)]$}
\RightLabel{\Elim{\lforall}}
\UnaryInfC{$!B(t)$}
\DisplayProof
\quad \longrightarrow\quad
\bottomAlignProof
\AxiomC{}
\RightLabel{$\Subst{\delta_2}{t}{c}$}
\shortDeduce
\DeduceC{$!B(t)$}
\DisplayProof}
\\[4ex]
\shoveright{\bottomAlignProof
\AxiomC{}
\RightLabel{$\delta_2$}
\shortDeduce
\DeduceC{$!B(t)$}
\RightLabel{\Intro{\lexists}}
\UnaryInfC{$\lexists[x][!B(x)]$}
\AxiomC{$\Discharge{!B(c)}{x}$}
\RightLabel{$\delta_3$}
\shortDeduce
\DeduceC{$!D$}
\DischargeRule{\Elim{\lexists}}{x}
\BinaryInfC{$!D$}
\DisplayProof
\quad \longrightarrow\quad
\bottomAlignProof
\AxiomC{}
\RightLabel{$\delta_2$}
\shortDeduce
\DeduceC{$!B(t)$}
\RightLabel{$\Subst{\Subst{\delta_3}{t}{c}}{\delta_2}{!B(t)^x}$}
\shortDeduce
\DeduceC{$!D$}
\DisplayProof}
\end{multline*}
\caption{Quelques conversions par réduction pour \Log{N1i}.}
\ollabel{tab:reduction-conversions}
\end{table}

\begin{prob}
Donnez la conversion par réduction pour
\Intro{\lnot} suivie de \Elim{\lnot}.
\end{prob}

Il reste le cas où l'occurrence de coupure~$!A$
provient de~\FalseInt. Il se traite de façon
analogue au cas d'une introduction.
Par exemple, on remplace
\[
\bottomAlignProof
\AxiomC{}
\RightLabel{$\delta_2$}
\DeduceC{$\lfalse$}
\RightLabel{\FalseInt}
\UnaryInfC{$!B \lif !C$}
\AxiomC{}
\RightLabel{$\delta_3$}
\DeduceC{$!B$}
\RightLabel{\Elim{\lif}}
\BinaryInfC{$!C$}
\DisplayProof
\quad\text{par}\quad
\bottomAlignProof
\AxiomC{}
\RightLabel{$\delta_2$}
\DeduceC{$\lfalse$}
\RightLabel{\FalseInt}
\UnaryInfC{$!C$}
\DisplayProof
\]
Une précaution supplémentaire est nécessaire
si $!A \ident (!B \lor !C)$ ou
$!A \ident \lexists[x][!B(x)]$.
Dans le cas disjonctif, le remplacement
\[
\bottomAlignProof
\AxiomC{}
\RightLabel{$\delta_2$}
\DeduceC{$\lfalse$}
\RightLabel{\FalseInt}
\UnaryInfC{$!B \lor !C$}
\AxiomC{$\Discharge{!B}{x}$}
\RightLabel{$\delta_3$}
\DeduceC{$!D$}
\AxiomC{$\Discharge{!C}{y}$}
\RightLabel{$\delta_4$}
\DeduceC{$!D$}
\DischargeRule{\Elim{\lor}}{x\,y}
\TrinaryInfC{$!D$}
\DisplayProof
\quad\text{par}\quad
\bottomAlignProof
\AxiomC{}
\RightLabel{$\delta_2$}
\DeduceC{$\lfalse$}
\RightLabel{\FalseInt}
\UnaryInfC{$!D$}
\DisplayProof
\]
pourrait créer une coupure sur~$!D$ ;
rien ne garantit
$\depth{!D} < \depth{!B \lor !C}$.
Il faut donc procéder comme dans la
réduction \Intro{\lor}/\Elim{\lor} :
on produit l'un des deux disjoints par
absurdité, puis on greffe sa preuve dans
la branche correspondante.
Notons les deux preuves prolongées de
l'absurde par le disjoint choisi à l'aide
d'un exposant sur le nom de la sous-preuve.
On obtient alors l'un des deux schémas :
\[
\bottomAlignProof
\AxiomC{}
\RightLabel{$\delta_2$}
\DeduceC{$\lfalse$}
\RightLabel{\FalseInt}
\UnaryInfC{$!B$}
\RightLabel{$\Subst{\delta_3}{\delta_2^{!B}}{!B^x}$}
\DeduceC{$!D$}
\DisplayProof
\quad\text{ou}\quad
\bottomAlignProof
\AxiomC{}
\RightLabel{$\delta_2$}
\DeduceC{$\lfalse$}
\RightLabel{\FalseInt}
\UnaryInfC{$!C$}
\RightLabel{$\Subst{\delta_4}{\delta_2^{!C}}{!C^y}$}
\DeduceC{$!D$}
\DisplayProof
\]
La notation de substitution indique le
greffage aux hypothèses désignées ; le
trait pointillé ne représente pas une
nouvelle règle allant arbitrairement
du disjoint à~$!D$.
Si la classe désignée est vide, la branche
mineure suffit à elle seule.

Dans le cas existentiel, on choisit un
terme fermé, on l'emploie pour instancier
la branche mineure à la place de la
constante propre, puis on dérive l'instance
de la matrice par~\FalseInt{} et on la
greffe aux hypothèses désignées.
Les seules coupures nouvelles au point
de greffage portent ainsi sur une instance
de la matrice, de rang inférieur à celui
de l'existentielle. Aux autres frontières,
les conditions de sélection excluent
l'allongement d'un segment de rang maximal.

















\end{document}
